\documentclass[12pt]{article}
\usepackage[utf8]{inputenc}
\usepackage[T1]{fontenc}
\usepackage{hyperref}
\title{Tapeba Language Resources for K-12 Curriculum Development}
\author{Compiled by Claude Sonnet 5 (claude-sonnet-5) for Andr\'as Kornai}
\date{}
\begin{document}
\maketitle

\section{Introduction}
Tapeba (ISO 639-3: tbb; Glottocode: tape1236) is the language of the Tapeba people of Caucaia, near Fortaleza in Cear\'a, Brazil. The search found nothing specific about the Tapeba language itself, such as a wordlist, grammar or documentation, and the Glottolog family label supplied in the prompt (``Bookkeeping'') carries no linguistic information. Search results were about Tapeba indigenous schooling, and they suggest the community now uses Portuguese. That is an inference from what the results did not show, and I did not confirm it.

\section{Primary Educational Resources}
\begin{itemize}
\item UECE journal article (\emph{CECiT}) on Tapeba differentiated schooling. It describes the Tapeba as pioneers of differentiated indigenous education in Cear\'a, with 14 such schools. \url{https://revistas.uece.br/index.php/CECiT/article/download/4734/8612/39493}
\item UFC undergraduate thesis (2023) on indigenous schooling, noting that the Escola Ind\'igena \'Indios Tapeba, in the early 1990s, was the first space for specific teaching of indigenous children in Cear\'a. \url{https://repositorio.ufc.br/bitstream/riufc/75430/1/2023_tcc_flsousa.pdf}
\item UFC book chapter on Tapeba education (2016). \url{https://repositorio.ufc.br/ri/bitstream/riufc/34491/1/2016_capliv_mfbfoliveiramhdieb.pdf}
\item \emph{Revista Research, Society and Development} article from the same search. I did not verify its content. \url{https://rsdjournal.org/index.php/rsd/article/download/24346/21447}
\end{itemize}
No Tapeba-language primers or textbooks were found.

\section{Linguistic and Reference Materials}
No dictionary, grammar or orthography for Tapeba was found.

\section{Religious and Community Resources}
\begin{itemize}
\item Di\'ario do Nordeste report on the Tapeba's roughly four-decade wait for territory demarcation in Caucaia (19 villages). \url{https://diariodonordeste.verdesmares.com.br/ceara/em-38-anos-de-luta-por-territorio-indigenas-tapeba-ainda-nao-conseguiram-demarcacao-em-caucaia-1.3340572}
\end{itemize}

\section{Digital and Audio Resources}
No audio, video or corpus resources for the Tapeba language were found.

\section{Curriculum Development Context}
The Tapeba already run differentiated indigenous schools, which is institutional groundwork for language work. The search did not show whether the ancestral language survives or what is recorded of it, so the first step would be a survey with community elders and archival research. Any curriculum would likely be Portuguese-medium with ancestral-language and culture content added as it is recovered.
\end{document}